\documentclass{situationist-preprint}

\title{Why the Efficacy of Prayer Cannot Be Measured}
\msid{REL-12-b}
\msfamily{REL-12}
\msdiscipline{Philosophy of Religion}
\mstype{Consequence}
\mskeywords{petitionary prayer; intercessory prayer trials; counterfactual dependence; granting; measurement}

\begin{document}
\maketitle

\begin{abstract}
REL-12-a argues that petitionary prayer can make a difference without causal efficacy: an outcome may be granted because it was requested even when it would otherwise have occurred. We derive why this efficacy cannot be measured. Outcome studies of intercessory prayer, from Galton's survey to randomized trials, compare the frequency of outcomes with and without prayer, and so test whether outcomes depend counterfactually on requests. The difference prayer makes on the account of REL-12-a is not a difference in outcomes but in whether an outcome is a granting. A null result is therefore consistent with every requested outcome that occurs being granted, and a positive result would establish causal efficacy, which is not what the account claims.
\end{abstract}

\section{What outcome studies test}

\citet{Galton1872} compared the longevity of those most prayed for, such as sovereigns, with that of other classes, and found no advantage. The randomized trials that followed, of which the largest is the STEP trial \citep{Benson2006}, assign patients to receive or not receive intercessory prayer and compare the frequency of a target outcome, such as freedom from complications, in the two groups.

\begin{proposition}\label{p:test}
A difference in outcome frequency between prayed-for and control groups, in a randomized design, is evidence that outcomes depend counterfactually on prayer; the absence of a difference is evidence that they do not.
\end{proposition}

\begin{proof}
Randomization makes the groups alike in expectation in everything except the prayer they receive. A difference in frequency is then evidence that, for some patients, the outcome would not have occurred had they not been prayed for, which is counterfactual dependence in the sense of REL-12-a's Definition~2.3; the absence of a difference is evidence against such dependence at the scale of the study.
\end{proof}

\section{What the account claims}

\begin{proposition}\label{p:invisible}
On the account of REL-12-a, every requested outcome that occurs may be granted, and yet the outcome frequencies in the two groups be the same.
\end{proposition}

\begin{proof}
By Lemma~3.1 of REL-12-a, an outcome may be granted in response to a request on which it does not counterfactually depend. Suppose that God brings about the same outcomes for patients in both groups, for reasons sufficient without the requests, and that for patients in the prayed-for group the requests are among the reasons for which God does so. Every requested outcome that occurs is then granted, by Definition~2.2 of REL-12-a, and the frequencies are equal.
\end{proof}

\begin{corollary}\label{c:measure}
No outcome study can measure the efficacy that REL-12-a attributes to petitionary prayer. A null result does not tell against it, and a positive result would establish a different thing, causal efficacy.
\end{corollary}

\begin{proof}
By Proposition~\ref{p:test}, outcome studies bear on counterfactual dependence of outcomes on prayer. By Proposition~\ref{p:invisible}, the presence or absence of granting is compatible with any frequencies, so it is not what such studies bear on. By Lemma~3.2 of REL-12-a, what depends on the request is the fact that the outcome is a granting, which is not an outcome that a study can count.
\end{proof}

The point is not that the relevant fact is hidden but that it is not of the kind outcome studies register. Whether an outcome is a granting is a fact about the reasons for which it was brought about. Reasons can be studied, in human agents, by asking and by observing patterns of choice across circumstances; but God's reasons, on the account, are not revealed by the pattern of outcomes, since by Principle~P1 of REL-12-a the same outcome may be brought about for several sufficient reasons at once. This is why, as REL-12-a puts it, answered and unanswered prayer can coincide in every observable respect.

\begin{objection}
A theory that no observation could disconfirm is empty.
\end{objection}

\begin{reply}
The account is not insulated from every observation. It is disconfirmed by evidence that there is no God, and it entails that petitioners can be heard, which is a claim about a relation between persons, not about frequencies. That a claim about the reasons for which an act is done cannot be tested by tabulating the act's effects is true of human gift-giving as well \citep{Stump1979}. What the corollary rules out is one method, the outcome study, for one question.
\end{reply}

\section{Conclusion}

Outcome studies test whether outcomes depend counterfactually on prayer. The efficacy of petitionary prayer, on the account of REL-12-a, lies in whether outcomes are granted, which does not depend on their frequency. The negative results of such studies, the STEP trial included, are consistent with the account, and no study of the same design could confirm it.

\begin{thebibliography}{99}\small
\bibitem[Benson et~al.(2006)]{Benson2006} Benson, H. et~al. (2006). Study of the therapeutic effects of intercessory prayer (STEP) in cardiac bypass patients: A multicenter randomized trial of uncertainty and certainty of receiving intercessory prayer. \emph{American Heart Journal}, 151(4), 934--942.
\bibitem[Galton(1872)]{Galton1872} Galton, F. (1872). Statistical inquiries into the efficacy of prayer. \emph{The Fortnightly Review}, 12, 125--135.
\bibitem[Stump(1979)]{Stump1979} Stump, E. (1979). Petitionary prayer. \emph{American Philosophical Quarterly}, 16(2), 81--91.
\end{thebibliography}

\end{document}
